\documentclass[11pt,a4paper]{article}
\usepackage{gtu-el-yazisi}
\begin{document}
\notbilgi{Alan Teorisi}{Elektrostatik potansiyel ve iş}{Ders6.pdf}

\sayfa{1}
\sekilnotu[1]{Düzensiz biçimli bir hacim ($v$); hacmin içinde orijinde buluşan $x'$ (sol-aşağı), $y'$ (sağ) ve $z'$ (yukarı) eksenleri; orijinden dışarıya, ``$\times$'' ile işaretli ve ``V'' etiketli bir noktaya giden $r'$ vektörü.}

\[
V=\frac{1}{4\pi\varepsilon_0}\int_{v}\frac{\rho}{r'}
=\frac{1}{4\pi\varepsilon_0}\iiint_{x',y',z'=-\infty}^{\infty}\frac{\rho(x',y',z')}{r'}\,dx'\,dy'\,dz'+V_0
\ \rightarrow\ \text{sabit}
\]

$\downarrow$ Coulomb potansiyeli \ $V(x,y,z)$

\begin{center}
$\displaystyle r'=\bigl[(x-x')+(y-y')+(z-z')\bigr]^{1/2}$\footnote{Claude notu: muhtemelen $r'=\bigl[(x-x')^2+(y-y')^2+(z-z')^2\bigr]^{1/2}$ olmalı.}
\end{center}

\[
\underbrace{\vec E}_{\text{p'deki elektrostatik alan}}=-\mathrm{grad}\,V
\]

\[
\mathrm{rot}\,\vec E=-\mathrm{rot}\,\mathrm{grad}\,V=0
\]

\[
\Delta V=\mathrm{div}\,\mathrm{grad}\,V=-\mathrm{div}\,\vec E
\qquad\qquad
\Delta=\vec\nabla\cdot\vec\nabla=\frac{\partial^2}{\partial x^2}+\frac{\partial^2}{\partial y^2}+\frac{\partial^2}{\partial z^2}
\]

\[
\downarrow\quad
=\frac{1}{4\pi\varepsilon_0}\int_{v}\rho(x',y',z')\,\Delta\!\left(\frac{1}{r'}\right)dv
\qquad r'\neq 0\ \text{ise}\quad \Delta\!\left(\frac{1}{r'}\right)=0\ \ (\rho\text{'nun olmadığı yerde})
\]

\[
\lim_{r\to\infty}V(x,y,z)\to 0 \qquad V_0=0\ \text{kabul edilir}
\]

\[
\underbrace{\Delta V=-\frac{\rho}{\varepsilon_0}}_{\text{Poisson denklemi}}\ \text{(her yerde)}
\qquad\qquad
\rho\triangleq\frac{dQ}{dv}
\]

\[
\begin{array}{rl}
V(x_1,y_1,z_1)&=\dfrac{1}{4\pi\varepsilon_0}\iiint\cdots+V_0\\
-\;V(x_2,y_2,z_2)&=\ ''\ \ \iiint\cdots+V_0\\ \hline
\end{array}
\]
\[
V_1-V_2=\frac{1}{4\pi\varepsilon_0}\left[\iiint\cdots-\iiint\cdots\right]+\underbrace{V_0-V_0}_{0}
\]

\sayfa{2}
\sekilnotu[2]{A noktasından B noktasına giden eğri bir yol (uçlar $\times$ ile işaretli); yol üzerinde bir noktada $q$ yükü ve ondan yukarı doğru çizilmiş $\vec F_q$ oku; yol boyunca teğet yönünde $d\vec c$ oku; A ve B yakınında kırmızı kalemle ``$q$'' yazılmış.}

\[
\vec F_q=q\cdot\vec E=-q\,\mathrm{grad}\,V=-q\,\vec\nabla\cdot V
\]

{\color{kirmizikalem}dışarıdan yapılan iş}
\[
\vec F=-\vec F_q=q\,\mathrm{grad}\,V
\]

\begin{align*}
W_{AB}&=\int_A^B\vec F\cdot d\vec c=\int_A^B q\,\mathrm{grad}\,V\,d\vec c\\
&=q\int_A^B\mathrm{grad}\,V\,d\vec c
\end{align*}

\[
dc\,\frac{\partial V}{\partial c}=\mathrm{grad}\,V\cdot\underbrace{\vec t\cdot dc}_{d\vec c}
\qquad\quad \vec t=d\vec\ell
\qquad\quad \frac{d\vec c}{dc}\triangleq\vec t
\]

\begin{align*}
&=q\int_A^B\mathrm{grad}\,V\cdot d\vec c\\
&=q\int_A^B\frac{\partial V}{\partial c}\cdot dc=q\,V\Big|_A^B\\
&=q\,V(B)-q\,V(A)=q\bigl[V(B)-V(A)\bigr]
\end{align*}

\[
q=1\,\mathrm{C}\ \text{(birim yük)}
\]
\[
A\to\infty\qquad V(A)\to 0\qquad\Longrightarrow\qquad =\ \underbrace{q}_{1\,\mathrm{C}}\,V(B)
\]

$V_B$: bir birim yük sonsuzdan sonlu bir B noktasına \okunamadi{getirinceye kadar} yapılan işe o noktadaki potansiyel denir.

$q=1\,\mathrm{C}$ yükü A'dan B'ye getirilinceye kadar yapılan işe potansiyel fark denir.

\sayfa{3}
\[
W_{AB}>0\ \Longrightarrow\ V_B>V_A
\]
\textcolor{kirmizikalem}{$\searrow$ Enerjiyi biz harcadık anlamına gelir\\ çünkü $\vec F$ bizim uyguladığımız kuvvet}

\medskip
\noindent$W_{AB}=0$ \quad ne zaman olur

\noindent\hspace*{1.2cm}$V_B=V(B_1)$

\noindent\hspace*{1.2cm}$V_A=V(A_1)$

\noindent\hspace*{1.2cm}$q=0$'sa olur

\noindent\hspace*{1.2cm}$q\neq 0$ için \quad $V_B=V_A$ olmalıdır

\noindent\hspace*{1.2cm}$A\to\infty$ \quad $W_{AB}=q\,V_B$

\[
\vec F_q=\frac{d}{dt}\bigl[m(v)\,\vec v\,\bigr]
\]
$\downarrow$ küçük hızlarda
\[
\vec F_q=m_0\cdot\frac{d\vec v}{dt}
\]
\[
d\vec c\cdot\vec F_q=m_0\cdot\frac{d\vec v}{dt}\cdot d\vec c
\]
\[
\vec F_q\cdot d\vec c=m_0\,\frac{d\vec v}{dt}\cdot\vec v\cdot dt
\qquad\qquad
\vec v\,/\!/\,d\vec v\ \Rightarrow\ \vec v\cdot d\vec v=v\cdot dv
\]

\begin{align*}
\int_A^B-\vec F\cdot d\vec c&=\int_A^B\vec F_q\cdot d\vec c\\
&=-q\int_A^B\mathrm{grad}\,V\,d\vec c=-q\,(V_B-V_A)\\
&=m_0\int_A^B\frac{d\vec v}{dt}\cdot\vec v\cdot dt=m_0\int_A^B v\,dv=\frac12\,m_0\,v^2\Big|_A^B\\
&=\frac12\,m_0\,(v_B^2-v_A^2)
\end{align*}

\formulkutu[Enerji korunur]{q\,V_B+\frac12\,m_0\,v_B^2=q\,V_A+\frac12\,m_0\,v_A^2}

\sayfa{4}
\textcolor{kirmizikalem}{$*$ İş yoldan bağımsızdır.}

\end{document}
